\subsection{Relations with functions of positive type}

\begin{enumerate}
\def\labelenumi{(\arabic{enumi})}
\setcounter{enumi}{65}
\tightlist
\item
\end{enumerate}

PROPOSITION 11. --- Let E be a finite-dimensional vector space, \(\mu\) a bounded (positive) measure on E, and \(\Phi\) the Fourier transform of \(\mu\) . The function \(\Phi\) is continuous and of positive type on \(E'\) .

The continuity of \(\Phi\) follows from Prop. 9 of No.~8.

Let us show that \(\Phi\) is of positive type. Let \(x_1', \ldots, x_p'\) be in E' and \(c_1, \ldots, c_p\) complex numbers. Then

\[
\begin{array}{r l} \sum_ {j, k} c _ {j} \overline {{{c}}} _ {k} \Phi (x _ {j} ^ {\prime} - x _ {k} ^ {\prime}) & = \int_ {\mathrm{E}} \sum_ {j, k} c _ {j} \overline {{{c}}} _ {k} e ^ {i \langle x, x _ {j} ^ {\prime} - x _ {k} ^ {\prime} \rangle} d \mu (x) \\ & = \int_ {\mathrm{E}} \left| \sum_ {j = 1} ^ {p} c _ {j} e ^ {i \langle x, x _ {j} ^ {\prime} \rangle} \right| ^ {2} d \mu (x) \geqslant 0. \end{array}\tag{Q.E.D.}
\]

One can prove a converse known as Bochner's theorem: every continuous function on \(E'\) of positive type is the Fourier transform of a bounded (positive) measure. \(^{*}\) We shall assume this result for the rest of No.~12.

THEOREM 4. --- Let E be a locally convex space. The Fourier transformation is a bijection of the set of promeasures on E onto the set of functions of positive type on E' whose restriction to every finite-dimensional subspace is continuous.

We know (No.~3, Prop. 3) that the Fourier transformation is injective. Let \(\mu = (\mu_{\mathrm{V}})_{\mathrm{V}\in \mathcal{F}(\mathbf{E})}\) be a promeasure on \(\mathbf{E}\) and \(\Phi\) its Fourier transform. Let \(\mathbf{T}\) be a finite-dimensional subspace of \(\mathbf{E}'\) and let \(\mathbf{V}\) be the orthogonal of \(\mathbf{T}\) in \(\mathbf{E}\) . One can identify \(\mathbf{T}\) with the dual of \(\mathbf{E} / \mathbf{V}\) ; the restriction \(\Phi_{\mathrm{T}}\) of \(\Phi\) to \(\mathbf{T}\) is the Fourier transform of the bounded measure \(\mu_{\mathrm{V}}\) on \(\mathbf{E} / \mathbf{V}\) . By Prop. 11, \(\Phi_{\mathrm{T}}\) is continuous and of positive type on \(\mathbf{T}\) . Since \(\mathbf{T}\) is arbitrary, it is clear that \(\Phi\) is of positive type on \(\mathbf{E}'\) .

Conversely, let \(\Phi\) be a function of positive type on \(\mathbf{E}'\) whose restriction to every finite-dimensional subspace of \(\mathbf{E}'\) is continuous. For every \(V \in \mathcal{F}(\mathbf{E})\) , we identify the dual of \(E/V\) with the orthogonal \(V^{\circ}\) of \(V\) in \(\mathbf{E}'\) ; the restriction \(\Phi_V\) of \(\Phi\) to \(V^{\circ}\) is continuous and of positive type and so, by Bochner's theorem, there exists a bounded (positive) measure \(\mu_V\) on \(E/V\) whose Fourier transform is \(\Phi_V\) . Let \(V\) and \(W\) in \(\mathcal{F}(\mathbf{E})\) be such that \(W \subset V\) , and let \(p_{\mathrm{VW}}\) be the canonical mapping of \(E/W\) onto \(E/V\) ; with the identifications made, \(^{t}p_{\mathrm{VW}}\) is the injection of \(V^{\circ}\) into \(W^{\circ}\) . By formula (4) of No.~3, we then have

\[
\mathcal {F} \big (p _ {\mathrm{VW}} (\mu_ {\mathrm{W}}) \big) = \big (\mathcal {F} \mu_ {\mathrm{W}} \big) \circ {} ^ {t} p _ {\mathrm{VW}} = \Phi_ {\mathrm{W}} \circ {} ^ {t} p _ {\mathrm{VW}} = \Phi_ {\mathrm{V}} = \mathcal {F} \mu_ {\mathrm{V}},
\]

whence \(p_{\mathrm{VW}}(\mu_{\mathrm{W}}) = \mu_{\mathrm{V}}\) by Prop. 3 of No.~3. Consequently, the family \(\mu = (\mu_{\mathrm{V}})_{\mathrm{V} \in \mathcal{F}(\mathrm{E})}\) is a promeasure on E; it is clear that \(\Phi\) is the Fourier transform of \(\mu\) .

COROLLARY. --- Let F be a barreled nuclear space; equip \(F'\) with a locally convex topology T intermediate to the weak topology \(\sigma(\mathbf{F}',\mathbf{F})\) and the topology of uniform convergence on the compact convex subsets of F. The Fourier transformation is a bijection of the set of bounded (positive) measures on \(F'\) onto the set of continuous functions of positive type on F.

This follows immediately from Theorem 4 and the Cor. of Th. 2 of No.~10.*

